% Figure: a rational function sketched from its features alone, with
% a vertical asymptote, a horizontal asymptote, a zero, and the sign
% behaviour across the pole. The point is that the engineer draws the
% curve from asymptotes and signs without plotting points.
\begin{figure}[htbp]
\centering
\begin{tikzpicture}[scale=0.95,
  ax/.style={->, draw=engblack, thick},
  tick/.style={draw=engblack!60, thin},
  asy/.style={draw=enggray, thin, dashed},
  curve/.style={draw=engblue, thick},
  ann/.style={font=\footnotesize, align=left},
]
% --- axes ---
\draw[ax] (-4.3, 0) -- (4.5, 0) node[below right, ann] {$x$};
\draw[ax] (0, -3.3) -- (0, 4.3) node[left, ann] {$y$};
\foreach \x in {-4, -3, -2, -1, 1, 2, 3, 4} {
  \draw[tick] (\x, -0.05) -- (\x, 0.05);
  \node[ann, below, yshift=-2pt] at (\x, 0) {$\x$};
}
\foreach \y in {-3, -2, -1, 1, 2, 3, 4} {
  \draw[tick] (-0.05, \y) -- (0.05, \y);
  \node[ann, left, xshift=-2pt] at (0, \y) {$\y$};
}
% r(x) = (x - 1)/(x + 2): vertical asymptote x=-2, horizontal y=1,
% zero at x=1.
% Vertical asymptote at x = -2
\draw[asy] (-2, -3.2) -- (-2, 4.2);
\node[ann, anchor=south west, color=enggray] at (-2, 4.0) {$x=-2$};
% Horizontal asymptote at y = 1
\draw[asy] (-4.2, 1) -- (4.4, 1);
\node[ann, anchor=south east, color=enggray] at (4.4, 1) {$y=1$};
% Left branch (x < -2): goes to +inf near asymptote, to 1 from above
% at -inf
\draw[curve, domain=-4.2:-2.25, samples=60, smooth]
  plot (\x, {(\x-1)/(\x+2)});
% Right branch (x > -2): comes from -inf, crosses zero at x=1,
% approaches 1 from below
\draw[curve, domain=-1.55:4.4, samples=80, smooth]
  plot (\x, {(\x-1)/(\x+2)});
% mark the zero
\node[circle, fill=engblue, inner sep=1.2pt] at (1, 0) {};
\node[ann, anchor=south west] at (1.05, 0.1) {zero $x=1$};
\node[ann, anchor=west, draw=engblue, fill=white, inner sep=2pt]
  at (1.8, 2.6) {$r(x) = \dfrac{x-1}{x+2}$};
\end{tikzpicture}
\caption[A rational function sketched from its features]{The rational
function $r(x) = (x - 1)/(x + 2)$ drawn from four features and no
plotted points. The vertical asymptote at $x = -2$ (the pole, where
the denominator vanishes) splits the plane into two branches. The
horizontal asymptote $y = 1$ (the ratio of leading coefficients) is
the limit at large $|x|$. The single zero at $x = 1$ is where the
numerator vanishes. The sign across the simple pole flips: the right
branch climbs from $-\infty$ just right of $x = -2$, through the zero
at $x = 1$, up toward $y = 1$ from below; the left branch falls from
$y = 1$ at the far left down toward $+\infty$ just left of the pole.
Four features fix the whole curve.}
\label{fig:vol02:ch02:rational-asymptotes}
\end{figure}
